\subsection{CONNECTED SETS IN A COMPACT SPACE}

DEFINITION 3. Let V be a symmetric entourage of a uniform space X. A finite sequence \((x_{i})_{0\leqslant i\leqslant n}\) of points of X is said to be a V-chain if \(x_{i}\) and \(x_{i+1}\) are V-close for \(0\leqslant i<n\) . The points \(x_{0}\) and \(x_{n}\) are called the ends of the V-chain, and they are said to be joined by the V-chain.

Given a symmetric entourage V, the relation ``there is a V-chain joining x and y'' is an equivalence relation between x and y in X. Let \(A_{x,v}\) be the equivalence class of x for this relation, i.e.~the set of all \(y \in X\) which can be joined to x by a V-chain. Clearly if \(y \in A_{x,v}\) then \(\mathrm{V}(y) \subset \mathrm{A}_{x,\mathrm{v}}\) ; hence \(A_{x,v}\) is open in X; and the complement of \(A_{x,v}\) , being a union of equivalence classes, is also open. Hence:

PROPOSITION 5. In a uniform space X, the set \(A_{x,v}\) of points which can be joined by a V-chain to a given point x is both open and closed in X.

For each \(x \in X\) , let \(A_{x}\) denote the intersection of the sets \(A_{x,v}\) as V runs through the set of symmetric entourages of X; \(A_{x}\) is the equivalence class of x for the equivalence relation ``for all symmetric entourages V, there is a V-chain joining x and y''.

PROPOSITION 6. In a compact space X, the component of x, the set \(A_{x}\) , and the intersection of the neighbourhoods of x which are both open and closed, are all three identical.

It is enough to show that \(A_{x}\) is connected: for in any uniform space X the component of x is contained in the intersection of the neighbourhoods of x which are both open and closed, and this intersection is contained in \(A_{x}\) by Proposition 5.

Suppose \(A_{x}\) is not connected. Since \(A_{x}\) is closed, there are two non-empty disjoint closed sets B and C such that \(B \cup C = A_{x}\) . By Proposition 4 of §3 there is an entourage U of X such that \(U(B) \cap U(C)\) is empty.

Let W be an open entourage such that \(\hat{W} \subset U\) , and let H be the closed set which is the complement of \(W(B) \cup W(C)\) in X. Suppose for example that \(x \in B\) , and let y be a point of C. Then for each symmetric entourage \(V \subset W\) one sees immediately, by induction on i, that every V-chain \((x_{i})_{0 \leqslant i \leqslant n}\) joining x and y in X must have a point in H, by the choice of W. Since by hypothesis x and y can be joined by a V-chain for each symmetric entourage V, we see that \(H \cap A_{x,v}\) is not empty if \(V \subset W\) . On the other hand, if \(V' \subset V\) then clearly \(A_{x,v'} \subset A_{x,v}\) ; thus it follows that, as V runs through the set of symmetric entourages of X, the sets \(H \cap A_{x,v}\) form a filter base of closed sets in the compact space H. Hence all these sets have a common point and therefore H meets \(A_{x}\) ; but this contradicts the definition of H.

COROLLARY. Let X be a locally compact space and let K be a compact component of X. Then the neighbourhoods of K which are both open and closed form a fundamental system of neighbourhoods of K.

Let V be a relatively compact open neighbourhood of K in X (Chapter I, § 9, no. 7, Proposition 10) and let F be its frontier. Let U ⊂ V be a set which is both open and closed with respect to V. Then U is closed in X, and if in addition U does not meet F, then U is open in X (for then U ⊂ V and U is open in V). Hence it is enough to show that there is a subset of V which contains K, does not meet F and is both open and closed with respect to V.

Suppose this is not the case: then the intersections of F with the subsets of \(\overline{V}\) which contain K and are open and closed in \(\overline{V}\) form a filter base of closed sets in F; F is compact, so these sets have a common point \(y \in F\) . But this is absurd; since \(\overline{V}\) is compact, K is a component of \(\overline{V}\) , and by virtue of Proposition 6, the intersection of the sets which contain K and are both open and closed in \(\overline{V}\) is just K. The Corollary is thus proved.

PROPOSITION 7. Let X be a compact space and let R be the equivalence relation on X whose classes are the components of X. Then the quotient space X/R is compact and totally disconnected.

We know from Chapter I, § II, no. 5, Proposition 9 that X/R is totally disconnected; hence we have to show that X/R is Hausdorff (Chapter I, § 10, no. 4, Proposition 8.) Let A and B be two distinct components of X. By Proposition 6 there is a symmetric entourage U of X such that no point of A can be joined to any point of B by a U-chain. The set V (resp. W) of points of X which can be joined to a point \(x \in A\) (resp. \(y \in B\) ) by a U-chain is both open and closed in X (Proposition 5) and contains A (resp. B); these sets are therefore open neighbourhoods of A and B respectively, are saturated with respect to R and do not intersect. This completes the proof.

